\chapter{Phát triển và triển khai ứng dụng}
\label{ch:dev}

\section{Thiết kế kiến trúc}

\subsection{Lựa chọn kiến trúc phần mềm}

Hệ thống \brandName{} được xây dựng theo kiến trúc phân tầng trên nền tảng Next.js, chia thành
ba tầng rõ ràng như Hình~\ref{fig:architecture}:

\begin{itemize}
  \item \textbf{Tầng trình diễn (Presentation):} chạy trên trình duyệt, gồm các thành phần giao
  diện React (Client Components) sử dụng Tailwind CSS để tạo kiểu và Zustand để quản lý trạng
  thái cục bộ như giỏ hàng, danh sách yêu thích và so sánh.
  \item \textbf{Tầng ứng dụng / nghiệp vụ (Application):} chạy trên máy chủ Next.js triển khai tại
  Vercel. Tầng này gồm các React Server Components (kết xuất giao diện và truy vấn dữ liệu đọc),
  Server Actions (xử lý các nghiệp vụ ghi như đặt hàng, kiểm duyệt), các Route Handler (webhook
  PayOS, tải ảnh, tìm kiếm) và tác vụ định kỳ (hủy đơn quá hạn).
  \item \textbf{Tầng dữ liệu \& dịch vụ (Data \& Services):} gồm cơ sở dữ liệu PostgreSQL trên
  Supabase (truy cập qua Prisma), dịch vụ xác thực và lưu trữ của Supabase, cùng các dịch vụ
  ngoài: PayOS (thanh toán), Resend (email) và Upstash (giới hạn truy cập).
\end{itemize}

\fig{figures/diagrams/architecture.pdf}{0.95}{\caption{Sơ đồ kiến trúc hệ thống ba tầng}\label{fig:architecture}}

Việc lựa chọn kiến trúc này tận dụng được thế mạnh của Next.js: phần lớn logic nghiệp vụ và truy
vấn dữ liệu được thực thi trên máy chủ, giúp giảm mã chạy phía client, tăng cường an toàn (không
lộ khóa bí mật ra trình duyệt) và tối ưu SEO --- yếu tố quan trọng đối với một website thương mại
điện tử.

\section{Thiết kế tổng quan}

\subsection{Mô tả tổng quát}

Hệ thống được tổ chức thành hai khu vực chính: khu vực bán hàng (storefront) dành cho khách
hàng và khu vực quản trị (admin) dành cho người bán. Mã nguồn được tổ chức theo tầng:

\begin{itemize}
  \item Tầng giao diện: các trang và thành phần trong thư mục \texttt{app/} và \texttt{components/}.
  \item Tầng nghiệp vụ: các Server Action và hàm dịch vụ trong \texttt{lib/server/}.
  \item Tầng truy vấn dữ liệu: các hàm đọc dữ liệu trong \texttt{lib/queries/}.
  \item Tầng kiểm tra dữ liệu: các lược đồ Zod trong \texttt{lib/validations/}.
\end{itemize}

\subsection{Luồng xử lý tổng quát}

Một yêu cầu đọc dữ liệu (ví dụ xem danh sách sản phẩm) được xử lý như sau: trình duyệt gửi yêu
cầu, máy chủ Next.js kết xuất Server Component, truy vấn dữ liệu qua Prisma, sau đó trả về HTML
đã được dựng sẵn cho trình duyệt. Với một thao tác ghi (ví dụ đặt hàng), thành phần giao diện gọi
trực tiếp một Server Action; hàm này kiểm tra dữ liệu bằng Zod, thực hiện nghiệp vụ trong một
giao dịch Prisma (kiểm tra tồn kho, tạo đơn, trừ kho, áp dụng voucher), rồi trả kết quả về giao
diện. Các sự kiện bất đồng bộ như thông báo thanh toán từ PayOS được xử lý qua Route Handler
webhook.

Hình~\ref{fig:seq-order} mô tả chi tiết luồng tương tác của nghiệp vụ đặt hàng và thanh toán ---
nghiệp vụ phức tạp và quan trọng nhất của hệ thống --- dưới dạng biểu đồ tuần tự, bao gồm cả hai
hình thức thanh toán COD và PayOS.

\fig{figures/diagrams/seq-order.pdf}{1.0}{\caption{Biểu đồ tuần tự nghiệp vụ đặt hàng và thanh toán}\label{fig:seq-order}}

\section{Thiết kế chi tiết}

\subsection{Thiết kế cơ sở dữ liệu}

Cơ sở dữ liệu của hệ thống được thiết kế gồm 17 bảng chính, mô tả bằng lược đồ Prisma và triển
khai trên PostgreSQL. Hình~\ref{fig:er} thể hiện mối liên hệ giữa các bảng. Các bảng chi tiết được
mô tả trong các Bảng~\ref{tab:db-user}--\ref{tab:db-banner}.

\fig{figures/diagrams/er.pdf}{1.0}{\caption{Mối liên hệ giữa các bảng trong cơ sở dữ liệu}\label{fig:er}}

\begin{dbtable}{Bảng User --- người dùng}\label{tab:db-user}
1 & id & uuid & PK & Định danh người dùng \\ \hline
2 & email & text & not null, unique & Địa chỉ email \\ \hline
3 & name & text & & Họ tên \\ \hline
4 & phone & text & & Số điện thoại \\ \hline
5 & role & enum & not null & Vai trò (USER/ADMIN) \\ \hline
6 & avatarUrl & text & & Ảnh đại diện \\ \hline
7 & createdAt & timestamp & not null & Ngày tạo \\ \hline
8 & updatedAt & timestamp & not null & Ngày cập nhật \\ \hline
\end{dbtable}

\begin{dbtable}{Bảng Address --- địa chỉ giao hàng}\label{tab:db-address}
1 & id & text & PK & Định danh địa chỉ \\ \hline
2 & userId & uuid & FK & Tham chiếu User \\ \hline
3 & recipientName & text & not null & Tên người nhận \\ \hline
4 & phone & text & not null & Số điện thoại \\ \hline
5 & province & text & not null & Tỉnh/thành \\ \hline
6 & district & text & not null & Quận/huyện \\ \hline
7 & ward & text & not null & Phường/xã \\ \hline
8 & addressLine & text & not null & Địa chỉ chi tiết \\ \hline
9 & isDefault & boolean & not null & Là địa chỉ mặc định \\ \hline
\end{dbtable}

\begin{dbtable}{Bảng Category --- danh mục sản phẩm}\label{tab:db-category}
1 & id & text & PK & Định danh danh mục \\ \hline
2 & name & text & not null & Tên danh mục \\ \hline
3 & slug & text & not null, unique & Đường dẫn thân thiện \\ \hline
4 & parentId & text & FK & Danh mục cha (tự tham chiếu) \\ \hline
5 & imageUrl & text & & Ảnh danh mục \\ \hline
6 & sortOrder & int & not null & Thứ tự hiển thị \\ \hline
\end{dbtable}

\begin{dbtable}{Bảng Product --- sản phẩm}\label{tab:db-product}
1 & id & text & PK & Định danh sản phẩm \\ \hline
2 & slug & text & not null, unique & Đường dẫn thân thiện \\ \hline
3 & name & text & not null & Tên sản phẩm \\ \hline
4 & description & text & not null & Mô tả \\ \hline
5 & categoryId & text & FK & Danh mục \\ \hline
6 & ageRange & enum[] & & Nhóm độ tuổi \\ \hline
7 & gender & enum & not null & Giới tính (BOY/GIRL/UNISEX) \\ \hline
8 & basePrice & int & not null & Giá gốc (VND) \\ \hline
9 & comparePrice & int & & Giá so sánh (gạch ngang) \\ \hline
10 & status & enum & not null & Trạng thái (DRAFT/ACTIVE/ARCHIVED) \\ \hline
11 & ratingAvg & float & not null & Điểm đánh giá trung bình \\ \hline
12 & ratingCount & int & not null & Số lượt đánh giá \\ \hline
\end{dbtable}

\begin{dbtable}{Bảng ProductImage --- ảnh sản phẩm}\label{tab:db-image}
1 & id & text & PK & Định danh ảnh \\ \hline
2 & productId & text & FK & Sản phẩm \\ \hline
3 & url & text & not null & Đường dẫn ảnh \\ \hline
4 & alt & text & not null & Văn bản thay thế \\ \hline
5 & sortOrder & int & not null & Thứ tự \\ \hline
6 & isPrimary & boolean & not null & Là ảnh đại diện \\ \hline
\end{dbtable}

\begin{dbtable}{Bảng ProductVariant --- biến thể sản phẩm}\label{tab:db-variant}
1 & id & text & PK & Định danh biến thể \\ \hline
2 & productId & text & FK & Sản phẩm \\ \hline
3 & sku & text & not null, unique & Mã SKU \\ \hline
4 & size & text & not null & Kích cỡ \\ \hline
5 & color & text & not null & Màu sắc \\ \hline
6 & colorHex & text & & Mã màu \\ \hline
7 & price & int & & Giá riêng của biến thể \\ \hline
8 & stock & int & not null & Số lượng tồn kho \\ \hline
\end{dbtable}

\begin{dbtable}{Bảng Order --- đơn hàng}\label{tab:db-order}
1 & id & text & PK & Định danh đơn hàng \\ \hline
2 & orderCode & text & not null, unique & Mã đơn hàng \\ \hline
3 & userId & uuid & FK & Khách hàng \\ \hline
4 & status & enum & not null & Trạng thái đơn \\ \hline
5 & paymentMethod & enum & not null & Phương thức thanh toán \\ \hline
6 & paymentStatus & enum & not null & Trạng thái thanh toán \\ \hline
7 & subtotal & int & not null & Tạm tính \\ \hline
8 & shippingFee & int & not null & Phí vận chuyển \\ \hline
9 & discountTotal & int & not null & Tổng giảm giá \\ \hline
10 & voucherCode & text & & Mã giảm giá áp dụng \\ \hline
11 & total & int & not null & Tổng thanh toán \\ \hline
12 & trackingCode & text & & Mã vận đơn \\ \hline
13 & shippingAddress & jsonb & not null & Địa chỉ giao (bản sao) \\ \hline
14 & payosOrderCode & bigint & unique & Mã đơn PayOS \\ \hline
15 & paymentExpiresAt & timestamp & & Hạn thanh toán \\ \hline
\end{dbtable}

\begin{dbtable}{Bảng OrderItem --- chi tiết đơn hàng}\label{tab:db-orderitem}
1 & id & text & PK & Định danh dòng hàng \\ \hline
2 & orderId & text & FK & Đơn hàng \\ \hline
3 & variantId & text & FK & Biến thể sản phẩm \\ \hline
4 & productName & text & not null & Tên sản phẩm (bản sao) \\ \hline
5 & variantSize & text & not null & Kích cỡ \\ \hline
6 & variantColor & text & not null & Màu sắc \\ \hline
7 & unitPrice & int & not null & Đơn giá \\ \hline
8 & quantity & int & not null & Số lượng \\ \hline
9 & subtotal & int & not null & Thành tiền \\ \hline
\end{dbtable}

\begin{dbtable}{Bảng InventoryLog --- nhật ký tồn kho}\label{tab:db-invlog}
1 & id & text & PK & Định danh bản ghi \\ \hline
2 & variantId & text & FK & Biến thể sản phẩm \\ \hline
3 & delta & int & not null & Lượng thay đổi (+/-) \\ \hline
4 & reason & enum & not null & Lý do thay đổi \\ \hline
5 & note & text & & Ghi chú \\ \hline
6 & adminUserId & uuid & FK & Quản trị viên thực hiện \\ \hline
7 & createdAt & timestamp & not null & Thời điểm \\ \hline
\end{dbtable}

\begin{dbtable}{Bảng Review --- đánh giá sản phẩm}\label{tab:db-review}
1 & id & text & PK & Định danh đánh giá \\ \hline
2 & productId & text & FK & Sản phẩm \\ \hline
3 & userId & uuid & FK & Người đánh giá \\ \hline
4 & rating & int & not null & Số sao (1--5) \\ \hline
5 & comment & text & not null & Nội dung nhận xét \\ \hline
6 & images & text[] & & Ảnh đính kèm \\ \hline
7 & status & enum & not null & Trạng thái kiểm duyệt \\ \hline
8 & adminReply & text & & Phản hồi của quản trị \\ \hline
\end{dbtable}

\begin{dbtable}{Bảng WishlistItem --- danh sách yêu thích}\label{tab:db-wishlist}
1 & id & text & PK & Định danh \\ \hline
2 & userId & uuid & FK & Người dùng \\ \hline
3 & productId & text & FK & Sản phẩm \\ \hline
4 & createdAt & timestamp & not null & Thời điểm thêm \\ \hline
\end{dbtable}

\begin{dbtable}{Bảng Voucher --- mã giảm giá}\label{tab:db-voucher}
1 & id & text & PK & Định danh voucher \\ \hline
2 & code & text & not null, unique & Mã giảm giá \\ \hline
3 & discountType & enum & not null & Loại giảm (PERCENT/FIXED) \\ \hline
4 & discountValue & int & not null & Giá trị giảm \\ \hline
5 & minOrderValue & int & not null & Giá trị đơn tối thiểu \\ \hline
6 & maxDiscount & int & & Mức giảm tối đa \\ \hline
7 & usageLimit & int & & Giới hạn lượt dùng \\ \hline
8 & perUserLimit & int & & Giới hạn theo người dùng \\ \hline
9 & startsAt & timestamp & not null & Thời điểm bắt đầu \\ \hline
10 & endsAt & timestamp & not null & Thời điểm kết thúc \\ \hline
\end{dbtable}

\begin{dbtable}{Bảng VoucherRedemption --- lượt dùng voucher}\label{tab:db-vr}
1 & id & text & PK & Định danh \\ \hline
2 & voucherId & text & FK & Voucher \\ \hline
3 & userId & uuid & FK & Người dùng \\ \hline
4 & orderId & text & FK, unique & Đơn hàng áp dụng \\ \hline
5 & discount & int & not null & Số tiền đã giảm \\ \hline
\end{dbtable}

\begin{dbtable}{Bảng FlashSale --- chương trình flash sale}\label{tab:db-flashsale}
1 & id & text & PK & Định danh chương trình \\ \hline
2 & name & text & not null & Tên chương trình \\ \hline
3 & startsAt & timestamp & not null & Thời điểm bắt đầu \\ \hline
4 & endsAt & timestamp & not null & Thời điểm kết thúc \\ \hline
5 & isActive & boolean & not null & Đang kích hoạt \\ \hline
\end{dbtable}

\begin{dbtable}{Bảng FlashSaleItem --- sản phẩm trong flash sale}\label{tab:db-fsi}
1 & id & text & PK & Định danh \\ \hline
2 & flashSaleId & text & FK & Chương trình flash sale \\ \hline
3 & productId & text & FK & Sản phẩm \\ \hline
4 & salePrice & int & not null & Giá khuyến mãi \\ \hline
\end{dbtable}

\begin{dbtable}{Bảng ReturnRequest --- yêu cầu đổi/trả}\label{tab:db-return}
1 & id & text & PK & Định danh yêu cầu \\ \hline
2 & orderId & text & FK, unique & Đơn hàng \\ \hline
3 & userId & uuid & FK & Khách hàng \\ \hline
4 & reason & text & not null & Lý do đổi/trả \\ \hline
5 & status & enum & not null & Trạng thái xử lý \\ \hline
6 & adminNote & text & & Ghi chú của quản trị \\ \hline
7 & resolvedAt & timestamp & & Thời điểm xử lý xong \\ \hline
\end{dbtable}

\begin{dbtable}{Bảng Banner --- banner trang chủ}\label{tab:db-banner}
1 & id & text & PK & Định danh banner \\ \hline
2 & title & text & not null & Tiêu đề \\ \hline
3 & imageUrl & text & not null & Ảnh banner \\ \hline
4 & linkUrl & text & & Liên kết khi nhấn \\ \hline
5 & sortOrder & int & not null & Thứ tự hiển thị \\ \hline
6 & isActive & boolean & not null & Đang hiển thị \\ \hline
\end{dbtable}

\subsection{Các kiểu liệt kê (enum)}

Bên cạnh các bảng, cơ sở dữ liệu sử dụng một số kiểu liệt kê (enum) để biểu diễn các giá trị có
tập hợp cố định, giúp ràng buộc dữ liệu chặt chẽ ngay ở tầng cơ sở dữ liệu. Bảng~\ref{tab:enums}
liệt kê các enum và giá trị tương ứng.

\begin{longtable}{|>{\raggedright\arraybackslash}p{0.26\linewidth}|>{\raggedright\arraybackslash}p{0.66\linewidth}|}
\caption{Các kiểu liệt kê (enum) trong cơ sở dữ liệu}\label{tab:enums}\\
\hline
\rowcolor{gray!12}\textbf{Tên enum} & \textbf{Các giá trị} \\ \hline
\endfirsthead
\hline
\rowcolor{gray!12}\textbf{Tên enum} & \textbf{Các giá trị} \\ \hline
\endhead
UserRole & USER, ADMIN \\ \hline
OrderStatus & PENDING, CONFIRMED, PACKING, SHIPPING, COMPLETED, CANCELED \\ \hline
PaymentMethod & COD, BANK\_TRANSFER, PAYOS \\ \hline
PaymentStatus & UNPAID, PAID, REFUNDED \\ \hline
ProductStatus & DRAFT, ACTIVE, ARCHIVED \\ \hline
Gender & BOY, GIRL, UNISEX \\ \hline
AgeRange & NEWBORN\_0\_1, TODDLER\_1\_3, KID\_3\_6, KID\_6\_12 \\ \hline
InventoryReason & MANUAL\_ADD, MANUAL\_REMOVE, ORDER\_RESERVE, ORDER\_CANCEL, STOCK\_TAKE \\ \hline
ReviewStatus & PENDING, APPROVED, REJECTED \\ \hline
DiscountType & PERCENT, FIXED \\ \hline
ReturnStatus & REQUESTED, APPROVED, REJECTED, REFUNDED \\ \hline
\end{longtable}

\subsection{Ràng buộc duy nhất và chỉ mục}

Để bảo đảm tính toàn vẹn và tối ưu hiệu năng truy vấn, cơ sở dữ liệu định nghĩa các ràng buộc
duy nhất (unique) và chỉ mục (index) trên những cột được truy vấn thường xuyên. Bảng~\ref{tab:indexes}
tổng hợp các ràng buộc và chỉ mục tiêu biểu.

\begin{longtable}{|>{\raggedright\arraybackslash}p{0.255\linewidth}|>{\raggedright\arraybackslash}p{0.32\linewidth}|>{\raggedright\arraybackslash}p{0.32\linewidth}|}
\caption{Ràng buộc duy nhất và chỉ mục tiêu biểu}\label{tab:indexes}\\
\hline
\rowcolor{gray!12}\textbf{Bảng} & \textbf{Ràng buộc duy nhất} & \textbf{Chỉ mục} \\ \hline
\endfirsthead
\hline
\rowcolor{gray!12}\textbf{Bảng} & \textbf{Ràng buộc duy nhất} & \textbf{Chỉ mục} \\ \hline
\endhead
User & email & --- \\ \hline
Category & slug & parentId, slug \\ \hline
Product & slug & (categoryId, status), (status, createdAt), slug \\ \hline
ProductVariant & sku & productId, sku \\ \hline
Order & orderCode, payosOrderCode & (userId, createdAt), (status, createdAt), (status, paymentStatus, paymentExpiresAt) \\ \hline
Review & (userId, productId) & (productId, status), (status, createdAt) \\ \hline
WishlistItem & (userId, productId) & (userId, createdAt) \\ \hline
Voucher & code & code, (isActive, startsAt, endsAt) \\ \hline
VoucherRedemption & orderId & voucherId, userId \\ \hline
FlashSaleItem & (flashSaleId, productId) & productId \\ \hline
ReturnRequest & orderId & (status, createdAt) \\ \hline
\end{longtable}

\section{Xây dựng ứng dụng}

\subsection{Thư viện và công cụ sử dụng}

Bảng~\ref{tab:libs} liệt kê các thư viện và công cụ chính được sử dụng trong quá trình xây dựng
hệ thống.

\begin{longtable}{|>{\raggedright\arraybackslash}p{0.30\linewidth}|>{\raggedright\arraybackslash}p{0.62\linewidth}|}
\caption{Danh sách thư viện và công cụ sử dụng}\label{tab:libs}\\
\hline
\rowcolor{gray!12}\textbf{Thư viện / Công cụ} & \textbf{Vai trò} \\ \hline
\endfirsthead
\hline
\rowcolor{gray!12}\textbf{Thư viện / Công cụ} & \textbf{Vai trò} \\ \hline
\endhead
Next.js 14 & Framework full-stack (App Router, RSC, Server Actions) \\ \hline
React 18 & Thư viện xây dựng giao diện \\ \hline
TypeScript & Ngôn ngữ lập trình với kiểu tĩnh \\ \hline
Prisma & ORM truy vấn cơ sở dữ liệu \\ \hline
PostgreSQL / Supabase & Cơ sở dữ liệu và hạ tầng backend \\ \hline
Tailwind CSS & Framework CSS utility-first \\ \hline
Zustand & Quản lý trạng thái phía client \\ \hline
React Hook Form + Zod & Quản lý và kiểm tra biểu mẫu \\ \hline
@payos/node & SDK tích hợp cổng thanh toán PayOS \\ \hline
Resend & Gửi email giao dịch \\ \hline
Upstash Redis & Giới hạn tần suất truy cập \\ \hline
SheetJS (xlsx) & Xuất báo cáo Excel \\ \hline
Recharts & Vẽ biểu đồ báo cáo \\ \hline
dnd-kit & Kéo thả sắp xếp banner \\ \hline
Vitest & Kiểm thử đơn vị \\ \hline
Playwright & Kiểm thử đầu cuối (E2E) \\ \hline
Vercel & Nền tảng triển khai \\ \hline
\end{longtable}

\subsection{Kết quả đạt được}

Hệ thống đã được xây dựng hoàn chỉnh với đầy đủ các chức năng đặt ra cho cả khu vực bán hàng
và khu vực quản trị. Giao diện được thiết kế theo phong cách hiện đại, thân thiện với phụ huynh,
hiển thị tốt trên nhiều kích thước màn hình. Phần dưới đây trình bày một số giao diện chính của
hệ thống.

\subsection{Giao diện các chức năng chính}

\screenshot{figures/screenshots/home.png}{0.9}{Giao diện trang chủ}{fig:home}
\screenshot{figures/screenshots/login.png}{0.9}{Giao diện trang đăng nhập}{fig:login}
\screenshot{figures/screenshots/register.png}{0.9}{Giao diện trang đăng ký}{fig:register}
\screenshot{figures/screenshots/products.png}{0.9}{Giao diện trang danh sách sản phẩm}{fig:products}
\screenshot{figures/screenshots/product-detail.png}{0.9}{Giao diện trang chi tiết sản phẩm}{fig:product-detail}
\screenshot{figures/screenshots/cart.png}{0.9}{Giao diện trang giỏ hàng}{fig:cart}
\screenshot{figures/screenshots/checkout.png}{0.9}{Giao diện trang thanh toán}{fig:checkout}
\screenshot{figures/screenshots/account-orders.png}{0.9}{Giao diện trang đơn hàng của khách hàng}{fig:account-orders}
\screenshot{figures/screenshots/account-wishlist.png}{0.9}{Giao diện trang danh sách yêu thích}{fig:account-wishlist}

Khu vực quản trị cung cấp các trang quản lý nghiệp vụ cho người bán. Một số giao diện quản trị
chính được trình bày dưới đây.

\screenshot{figures/screenshots/admin-products.png}{0.9}{Giao diện trang quản trị sản phẩm}{fig:admin-products}
\screenshot{figures/screenshots/admin-inventory.png}{0.9}{Giao diện trang quản trị tồn kho}{fig:admin-inventory}
\screenshot{figures/screenshots/admin-orders.png}{0.9}{Giao diện trang quản trị đơn hàng}{fig:admin-orders}
\screenshot{figures/screenshots/admin-customers.png}{0.9}{Giao diện trang quản trị khách hàng}{fig:admin-customers}
\screenshot{figures/screenshots/admin-vouchers.png}{0.9}{Giao diện trang quản lý mã giảm giá}{fig:admin-vouchers}
\screenshot{figures/screenshots/admin-flash-sales.png}{0.9}{Giao diện trang quản lý flash sale}{fig:admin-flash-sales}
\screenshot{figures/screenshots/admin-banners.png}{0.9}{Giao diện trang quản lý banner trang chủ}{fig:admin-banners}
\screenshot{figures/screenshots/admin-reports.png}{0.9}{Giao diện trang thống kê doanh thu}{fig:admin-reports}

\section{Kiểm thử}

\subsection{Xây dựng một số testcase và kết quả đạt được}

Hệ thống được kiểm thử ở hai mức: kiểm thử đơn vị (unit test) bằng Vitest cho các hàm nghiệp vụ
thuần (tính giá, áp dụng voucher, xử lý thanh toán PayOS) và kiểm thử đầu cuối bằng Playwright
cho các luồng người dùng quan trọng. Bảng~\ref{tab:testcases} trình bày một số ca kiểm thử tiêu
biểu và kết quả.

\begin{longtable}{|c|>{\raggedright\arraybackslash}p{0.30\linewidth}|>{\raggedright\arraybackslash}p{0.30\linewidth}|c|}
\caption{Một số ca kiểm thử chức năng}\label{tab:testcases}\\
\hline
\rowcolor{gray!12}\textbf{Mã} & \textbf{Nội dung kiểm thử} & \textbf{Kết quả mong đợi} & \textbf{Trạng thái} \\ \hline
\endfirsthead
\hline
\rowcolor{gray!12}\textbf{Mã} & \textbf{Nội dung kiểm thử} & \textbf{Kết quả mong đợi} & \textbf{Trạng thái} \\ \hline
\endhead
TC01 & Tính giá hiệu lực khi có flash sale & Trả về giá khuyến mãi & Đạt \\ \hline
TC02 & Áp dụng voucher giảm theo phần trăm & Giảm đúng số tiền, không vượt mức tối đa & Đạt \\ \hline
TC03 & Áp dụng voucher hết hạn & Từ chối, báo lỗi & Đạt \\ \hline
TC04 & Đặt hàng khi đủ tồn kho & Tạo đơn, trừ tồn kho & Đạt \\ \hline
TC05 & Đặt hàng khi sản phẩm hết hàng & Hủy giao dịch, báo lỗi & Đạt \\ \hline
TC06 & Xác nhận thanh toán PayOS hợp lệ & Cập nhật trạng thái đã thanh toán & Đạt \\ \hline
TC07 & Webhook PayOS gọi lặp lại & Idempotent, không xử lý trùng & Đạt \\ \hline
TC08 & Hủy đơn PayOS quá hạn & Hủy đơn, hoàn trả tồn kho & Đạt \\ \hline
TC09 & Đăng nhập sai mật khẩu nhiều lần & Bị giới hạn truy cập & Đạt \\ \hline
TC10 & Kiểm duyệt đánh giá & Cập nhật điểm trung bình sản phẩm & Đạt \\ \hline
\end{longtable}

\section{Triển khai}

Hệ thống được triển khai trên nền tảng Vercel theo quy trình tích hợp và triển khai liên tục: mỗi
khi mã nguồn được đẩy lên kho, Vercel tự động build và triển khai phiên bản mới. Cơ sở dữ liệu
PostgreSQL được vận hành trên Supabase. Các biến môi trường (khóa kết nối cơ sở dữ liệu, khóa
PayOS, Resend, Upstash) được cấu hình an toàn trong bảng điều khiển của Vercel. Một tác vụ cron
được cấu hình chạy định kỳ (mỗi 15 phút) để quét và hủy các đơn hàng PayOS chưa thanh toán
quá hạn, đồng thời hoàn trả tồn kho. Bảng~\ref{tab:deploy} tóm tắt môi trường triển khai.

\begin{longtable}{|>{\raggedright\arraybackslash}p{0.32\linewidth}|>{\raggedright\arraybackslash}p{0.60\linewidth}|}
\caption{Môi trường triển khai hệ thống}\label{tab:deploy}\\
\hline
\rowcolor{gray!12}\textbf{Thành phần} & \textbf{Mô tả} \\ \hline
\endfirsthead
\hline
\rowcolor{gray!12}\textbf{Thành phần} & \textbf{Mô tả} \\ \hline
\endhead
Ứng dụng web & Next.js triển khai trên Vercel (serverless) \\ \hline
Cơ sở dữ liệu & PostgreSQL trên Supabase \\ \hline
Lưu trữ ảnh & Supabase Storage \\ \hline
Tác vụ định kỳ & Vercel Cron (hủy đơn quá hạn, mỗi 15 phút) \\ \hline
Tên miền & Cấu hình qua Vercel \\ \hline
\end{longtable}
